\chapter{Accessibility and Assistive Modes}\label{record-032}

\noindent{\footnotesize\color{BookMuted}\path{docs/strategy/accessibility-and-assistive-modes.md}\par}\smallskip

\emph{Design note for accessible Bijou experiences beyond visual rich-mode parity}

\section{Why this exists}\label{record-032-why-this-exists}

Bijou already has an \texttt{a\allowbreak{}c\allowbreak{}c\allowbreak{}e\allowbreak{}s\allowbreak{}s\allowbreak{}i\allowbreak{}b\allowbreak{}l\allowbreak{}e} lowering target for many individual components.

That is useful, but it is not enough.

A component can lower honestly while the overall app still fails people who need:

\begin{itemize}
\tightlist
\item
  screen readers
\item
  low-vision support
\item
  color-blind-safe signaling
\item
  reduced motion
\item
  lower cognitive load
\item
  alternate input or assistive interaction
\end{itemize}

This note exists to make one thing explicit:

\texttt{a\allowbreak{}c\allowbreak{}c\allowbreak{}e\allowbreak{}s\allowbreak{}s\allowbreak{}i\allowbreak{}b\allowbreak{}l\allowbreak{}e} output is necessary, but it is not the same thing as an accessible product.

\section{Current reality}\label{record-032-current-reality}

Today Bijou is relatively strong at:

\begin{itemize}
\tightlist
\item
  giving many core components an explicit \texttt{a\allowbreak{}c\allowbreak{}c\allowbreak{}e\allowbreak{}s\allowbreak{}s\allowbreak{}i\allowbreak{}b\allowbreak{}l\allowbreak{}e} rendering path
\item
  avoiding purely decorative borders/backgrounds in accessible mode
\item
  encouraging explicit text like \texttt{W\allowbreak{}a\allowbreak{}r\allowbreak{}n\allowbreak{}i\allowbreak{}n\allowbreak{}g\allowbreak{}:\allowbreak{} \allowbreak{}m\allowbreak{}e\allowbreak{}s\allowbreak{}s\allowbreak{}a\allowbreak{}g\allowbreak{}e} instead of ornamental styling
\end{itemize}

Today Bijou is not yet strong enough at:

\begin{itemize}
\tightlist
\item
  making a complex framed app like DOGFOOD truly usable in accessible mode
\item
  representing interactive structure in a semantic layer above pure visual layout
\item
  addressing low-vision and color-blind needs as first-class design inputs
\end{itemize}

\section{Product stance}\label{record-032-product-stance}

Bijou should treat accessibility as a design system concern, not a fallback string concern.

That means:

\begin{itemize}
\tightlist
\item
  a view may need to lower to a different experience, not merely a flatter rendering
\item
  linear reading order matters more than preserving spatial cleverness
\item
  visual semantics must become explicit text when required
\item
  assistive users should not be abandoned simply because a surface was originally designed for rich TUI interaction
\end{itemize}

\section{Design principles}\label{record-032-design-principles}

\subsection{1. Linear truth beats spatial fidelity}\label{record-032-linear-truth-beats-spatial-fidelity}

When an interface becomes hard to perceive spatially, Bijou should prefer a linear, explicit representation over a compressed imitation of the rich layout.

For example:

\begin{itemize}
\tightlist
\item
  a three-pane docs shell should probably lower to a guided sequential reading model
\item
  not a cramped textual copy of the same spatial arrangement
\end{itemize}

\subsection{2. Important meaning must survive without color}\label{record-032-important-meaning-must-survive-without-color}

State should not rely on hue alone.

Selection, severity, status, and affordance should also be carried by:

\begin{itemize}
\tightlist
\item
  text
\item
  labels
\item
  structure
\item
  glyphs
\item
  grouping
\end{itemize}

\subsection{3. Reduced motion is a first-class mode}\label{record-032-reduced-motion-is-a-first-class-mode}

Animation can enrich a TUI, but it should never be required for comprehension.

Bijou should support:

\begin{itemize}
\tightlist
\item
  reduced-motion profiles
\item
  low-motion shell behavior
\item
  motion-free explanations of dynamic state
\end{itemize}

\subsection{4. Accessible mode is not a dump}\label{record-032-accessible-mode-is-not-a-dump}

Plain text should still be:

\begin{itemize}
\tightlist
\item
  structured
\item
  oriented
\item
  navigable
\item
  intentional
\end{itemize}

Accessible lowering should not mean ``unstyled but still confusing.''

\subsection{5. Cognitive load is part of accessibility}\label{record-032-cognitive-load-is-part-of-accessibility}

Accessibility includes how much the user has to remember, infer, and hold in working memory.

We should minimize:

\begin{itemize}
\tightlist
\item
  hidden modes
\item
  overloaded hint bars
\item
  unexplained jargon
\item
  surfaces with too many equally weighted actions
\end{itemize}

\section{Assistive posture}\label{record-032-assistive-posture}

\subsection{Screen readers and terminal reality}\label{record-032-screen-readers-and-terminal-reality}

Terminal screen-reader behavior varies and may not provide a rich widget model comparable to the web.

Bijou should therefore not assume:

\begin{itemize}
\tightlist
\item
  semantic terminal accessibility metadata will rescue a visually dense TUI
\item
  a screen reader can infer pane relationships or control ownership from layout alone
\end{itemize}

Instead, Bijou should own semantic clarity in the content it emits.

\subsection{Voice and alternate interaction}\label{record-032-voice-and-alternate-interaction}

Bijou should eventually grow toward a semantic model that can describe:

\begin{itemize}
\tightlist
\item
  role
\item
  label
\item
  description
\item
  current state
\item
  available actions
\item
  focus order
\end{itemize}

That model could support:

\begin{itemize}
\tightlist
\item
  accessible lowering
\item
  docs
\item
  testing
\item
  explainability
\item
  future assistive adapters
\end{itemize}

This should be treated as a semantic shadow model rather than a visual clone of rich-mode layout.

\section{Low-vision design posture}\label{record-032-low-vision-design-posture}

Bijou should eventually support and audit for:

\begin{itemize}
\tightlist
\item
  high-contrast themes
\item
  monochrome-safe or low-hue themes
\item
  reduced visual clutter
\item
  larger spacing or simplified shell modes where appropriate
\end{itemize}

This does not mean every app must ship all of these immediately.

It does mean the framework should make them possible and should not fight them.

\section{Color-blindness posture}\label{record-032-color-blindness-posture}

The framework should help builders avoid:

\begin{itemize}
\tightlist
\item
  meaning carried only by red/green or similar hue pairings
\item
  state changes visible only through chroma
\item
  token combinations that collapse under common color-vision deficiencies
\end{itemize}

Future tooling should include:

\begin{itemize}
\tightlist
\item
  contrast checks
\item
  token-pair warnings
\item
  color-blind simulation or heuristics
\end{itemize}

\section{DOGFOOD implication}\label{record-032-dogfood-implication}

DOGFOOD is the proving surface, which means Bijou should eventually support an accessible DOGFOOD experience that is genuinely usable.

The likely direction is:

\begin{itemize}
\tightlist
\item
  search-first
\item
  linear navigation
\item
  explicit ``you are here'' copy
\item
  explicit action lists
\item
  no dependence on three-pane spatial inference
\end{itemize}

That should be treated as a different truthful experience, not a cramped copy of the rich shell.

That is a different experience, not just a visually flattened copy of the rich shell.

\section{Backlog direction}\label{record-032-backlog-direction}

Recommended next-layer backlog:

\begin{itemize}
\tightlist
\item
  high-contrast / low-motion shell profiles
\item
  token contrast and color-blind auditing tools
\item
  semantic shadow-model exploration for interactive surfaces
\item
  accessibility review of DOGFOOD as a full app, not just component stories
\item
  cognitive-load heuristics for shell and component surfaces
\end{itemize}

\section{Relationship to other doctrine}\label{record-032-relationship-to-other-doctrine}

This note should be read alongside:

\begin{itemize}
\tightlist
\item
  \href{https://github.com/flyingrobots/bijou/blob/cbd70612140e7a450d764b3c7be1c6af878b99bc/docs/strategy/bijou-ux-doctrine.md}{Bijou UX Doctrine}
\item
  \href{https://github.com/flyingrobots/bijou/blob/cbd70612140e7a450d764b3c7be1c6af878b99bc/docs/strategy/humane-shell.md}{The Humane Shell}
\item
  \href{https://github.com/flyingrobots/bijou/blob/cbd70612140e7a450d764b3c7be1c6af878b99bc/docs/strategy/localization-and-bidirectionality.md}{Localization and Bidirectionality}
\item
  \href{https://github.com/flyingrobots/bijou/blob/cbd70612140e7a450d764b3c7be1c6af878b99bc/docs/strategy/content-guide.md}{Content Guide}
\end{itemize}

\section{Acceptance bar}\label{record-032-acceptance-bar}

This direction is working when:

\begin{itemize}
\tightlist
\item
  Bijou stops equating accessible mode with ``flat text fallback''
\item
  rich-shell surfaces are evaluated as experiences, not just as render outputs
\item
  low-vision, motion, and color-blind needs influence design decisions early
\item
  the framework begins to grow explicit semantic and audit tooling in support of that goal
\end{itemize}

